% ========================================================
% BÀI TẬP TOÁN 5 - KẾT NỐI TRI THỨC VỚI CUỘC SỐNG
% Bài 13: Làm tròn số thập phân
% Định dạng: BTTeX Part III - Trắc nghiệm điền số ngắn
% ========================================================

% Câu 1
\begin{ex}
  Làm tròn số thập phân $4{,}7$ đến hàng đơn vị ta được số nào?
  \shortans{5}
  \loigiai{
    Chữ số hàng phần mười là 7 ($7 \ge 5$) nên làm tròn lên thành 5.
  }
\end{ex}

% Câu 2
\begin{ex}
  Làm tròn số thập phân $8{,}3$ đến hàng đơn vị ta được số nào?
  \shortans{8}
  \loigiai{
    Chữ số hàng phần mười là 3 ($3 < 5$) nên làm tròn xuống thành 8.
  }
\end{ex}

% Câu 3
\begin{ex}
  Làm tròn số thập phân $15{,}5$ đến hàng đơn vị ta được số nào?
  \shortans{16}
  \loigiai{
    Chữ số hàng phần mười là 5 ($5 \ge 5$) nên làm tròn lên thành 16.
  }
\end{ex}

% Câu 4
\begin{ex}
  Làm tròn số thập phân $24{,}38$ đến hàng phần mười ta được số $24{,}4$. Chữ số hàng phần mười của số sau khi làm tròn là bao nhiêu?
  \shortans{4}
  \loigiai{
    Chữ số hàng phần trăm là 8 ($8 \ge 5$) nên làm tròn lên được $24{,}4$. Chữ số hàng phần mười là 4.
  }
\end{ex}

% Câu 5
\begin{ex}
  Làm tròn số thập phân $9{,}72$ đến hàng phần mười ta được số $9{,}7$. Chữ số hàng phần mười là bao nhiêu?
  \shortans{7}
  \loigiai{
    Chữ số hàng phần trăm là 2 ($2 < 5$) nên làm tròn xuống được $9{,}7$. Chữ số là 7.
  }
\end{ex}

% Câu 6
\begin{ex}
  Làm tròn số thập phân $3{,}648$ đến hàng phần trăm ta được số $3{,}65$. Chữ số hàng phần trăm sau khi làm tròn là bao nhiêu?
  \shortans{5}
  \loigiai{
    Chữ số hàng phần nghìn là 8 ($8 \ge 5$) nên làm tròn lên thành $3{,}65$. Chữ số hàng phần trăm là 5.
  }
\end{ex}

% Câu 7
\begin{ex}
  Làm tròn số thập phân $7{,}824$ đến hàng phần trăm ta được số $7{,}82$. Chữ số hàng phần trăm sau khi làm tròn là bao nhiêu?
  \shortans{2}
  \loigiai{
    Chữ số hàng phần nghìn là 4 ($4 < 5$) nên làm tròn xuống thành $7{,}82$. Chữ số là 2.
  }
\end{ex}

% Câu 8
\begin{ex}
  Làm tròn số thập phân $99{,}6$ đến hàng đơn vị ta được số nào?
  \shortans{100}
  \loigiai{
    Chữ số hàng phần mười là 6 ($6 \ge 5$) nên làm tròn lên thành 100.
  }
\end{ex}

% Câu 9
\begin{ex}
  Có bao nhiêu số thập phân có một chữ số ở phần thập phân khi làm tròn đến hàng đơn vị thì được số $8$?
  \shortans{10}
  \loigiai{
    Các số từ $7{,}5$ đến $8{,}4$: $7{,}5; 7{,}6; 7{,}7; 7{,}8; 7{,}9; 8{,}0; 8{,}1; 8{,}2; 8{,}3; 8{,}4$. Có tất cả 10 số.
  }
\end{ex}

% Câu 10
\begin{ex}
  Một túi cam cân nặng $3{,}78\text{ kg}$. Làm tròn khối lượng túi cam đến hàng đơn vị ta được bao nhiêu ki-lô-gam?
  \shortans{4}
  \loigiai{
    Chữ số hàng phần mười là 7 nên làm tròn lên thành $4\text{ kg}$.
  }
\end{ex}

% Câu 11
\begin{ex}
  Chiều dài một sợi dây là $12{,}4\text{ m}$. Làm tròn độ dài đến hàng đơn vị ta được bao nhiêu mét?
  \shortans{12}
  \loigiai{
    Chữ số hàng phần mười là 4 nên làm tròn xuống thành $12\text{ m}$.
  }
\end{ex}

% Câu 12
\begin{ex}
  Làm tròn số $0{,}85$ đến hàng phần mười ta được số $0{,}9$. Chữ số hàng phần mười là bao nhiêu?
  \shortans{9}
  \loigiai{
    Chữ số hàng phần trăm là 5 nên làm tròn lên thành $0{,}9$. Chữ số là 9.
  }
\end{ex}

% Câu 13
\begin{ex}
  Làm tròn số $5{,}09$ đến hàng phần mười ta được số $5{,}1$. Chữ số hàng phần mười là bao nhiêu?
  \shortans{1}
  \loigiai{
    Chữ số hàng phần trăm là 9 nên làm tròn lên thành $5{,}1$. Chữ số là 1.
  }
\end{ex}

% Câu 14
\begin{ex}
  Làm tròn số $45{,}206$ đến hàng phần trăm ta được số $45{,}21$. Chữ số hàng phần trăm là bao nhiêu?
  \shortans{1}
  \loigiai{
    Chữ số hàng phần nghìn là 6 nên làm tròn lên thành $45{,}21$. Chữ số là 1.
  }
\end{ex}

% Câu 15
\begin{ex}
  Tìm chữ số $x$ biết rằng khi làm tròn số $14{,}x$ đến hàng đơn vị ta được số $14$. Có bao nhiêu chữ số $x$ thỏa mãn?
  \shortans{5}
  \loigiai{
    Để làm tròn xuống 14 thì $x \in \{0, 1, 2, 3, 4\}$. Có 5 chữ số.
  }
\end{ex}

% Câu 16
\begin{ex}
  Tìm chữ số $y$ biết rằng khi làm tròn số $25{,}y$ đến hàng đơn vị ta được số $26$. Có bao nhiêu chữ số $y$ thỏa mãn?
  \shortans{5}
  \loigiai{
    Để làm tròn lên 26 thì $y \in \{5, 6, 7, 8, 9\}$. Có 5 chữ số.
  }
\end{ex}

% Câu 17
\begin{ex}
  Một can dầu chứa $19{,}65$ lít dầu. Làm tròn thể tích dầu đến hàng đơn vị ta được bao nhiêu lít?
  \shortans{20}
  \loigiai{
    Chữ số hàng phần mười là 6 nên làm tròn lên thành 20 lít.
  }
\end{ex}

% Câu 18
\begin{ex}
  Số thập phân có một chữ số ở phần thập phân lớn nhất khi làm tròn đến hàng đơn vị được số $15$ là $15{,}4$. Chữ số hàng phần mười của số đó là bao nhiêu?
  \shortans{4}
  \loigiai{
    Số lớn nhất là $15{,}4$. Chữ số hàng phần mười là 4.
  }
\end{ex}

% Câu 19
\begin{ex}
  Số thập phân có một chữ số ở phần thập phân bé nhất khi làm tròn đến hàng đơn vị được số $15$ là $14{,}5$. Chữ số hàng phần mười của số đó là bao nhiêu?
  \shortans{5}
  \loigiai{
    Số bé nhất là $14{,}5$. Chữ số hàng phần mười là 5.
  }
\end{ex}

% Câu 20
\begin{ex}
  Làm tròn số $345{,}82$ đến hàng chục ta được số nào?
  \shortans{350}
  \loigiai{
    Chữ số hàng đơn vị là 5 nên làm tròn lên thành 350.
  }
\end{ex}

% Câu 21
\begin{ex}
  Làm tròn số $123{,}45$ đến hàng trăm ta được số nào?
  \shortans{100}
  \loigiai{
    Chữ số hàng chục là 2 ($2 < 5$) nên làm tròn xuống thành 100.
  }
\end{ex}

% Câu 22
\begin{ex}
  Tính tổng $4{,}6 + 3{,}8 = 8{,}4$. Làm tròn kết quả tổng này đến hàng đơn vị được số nào?
  \shortans{8}
  \loigiai{
    Tổng là $8{,}4$. Chữ số hàng phần mười là 4 nên làm tròn xuống thành 8.
  }
\end{ex}

% Câu 23
\begin{ex}
  Tính hiệu $12{,}5 - 4{,}2 = 8{,}3$. Làm tròn kết quả hiệu này đến hàng đơn vị được số nào?
  \shortans{8}
  \loigiai{
    Hiệu là $8{,}3$. Làm tròn đến hàng đơn vị được 8.
  }
\end{ex}

% Câu 24
\begin{ex}
  Tính tích $2{,}5 \times 3 = 7{,}5$. Làm tròn kết quả này đến hàng đơn vị được số nào?
  \shortans{8}
  \loigiai{
    Tích là $7{,}5$. Chữ số hàng phần mười là 5 nên làm tròn lên thành 8.
  }
\end{ex}

% Câu 25
\begin{ex}
  Một học sinh có điểm trung bình các bài kiểm tra là $8{,}65$. Điểm số này sau khi làm tròn đến hàng phần mười là $8{,}7$. Chữ số hàng phần mười là bao nhiêu?
  \shortans{7}
  \loigiai{
    Chữ số hàng phần trăm là 5 nên làm tròn lên thành $8{,}7$. Chữ số hàng phần mười là 7.
  }
\end{ex}
